% Generated draft: requires mathematical and source review.
\documentclass[11pt]{article}
\usepackage[margin=25mm]{geometry}
\usepackage{fontspec}
\setmainfont{Latin Modern Roman}
\usepackage{amsmath,amssymb,mathtools,unicode-math}
\setmathfont{Latin Modern Math}
\usepackage{enumitem,needspace,xurl,hyperref}
\setlength{\parindent}{0pt}
\setlength{\parskip}{5pt}
\setlength{\emergencystretch}{3em}
\newcommand{\N}{\mathbb{N}}
\newcommand{\Z}{\mathbb{Z}}
\newcommand{\Q}{\mathbb{Q}}
\newcommand{\R}{\mathbb{R}}
\newcommand{\C}{\mathbb{C}}
\newcommand{\F}{\mathbb{F}}
\newcommand{\A}{\mathbb{A}}
\newcommand{\PP}{\mathbb{P}}
\newcommand{\E}{\mathbb{E}}
\newcommand{\eps}{\varepsilon}
\newcommand{\dd}{\,\mathrm{d}}
\begin{document}
\section*{UnsolvedMath ID: 30005430}
\textbf{Source problem code:} catalogue entry
\subsection{Definitions and mathematical statement}
The supplied record does not determine a standalone open mathematical problem. Its title asks whether branching now implies branching next, but its statement-verification field identifies this as a talk heading introducing an already developed statistical test and marks the mathematical statement as unrecoverable. The following records the reported null-model relation and its unresolved qualifications.

For a branching node \(v\) in a timed phylogeny, let \(m_v\) be its number of child lineages and let \(k_v\) be the total number of lineages counted at that node's time. Let \(A_v\) denote the event that one of those child lineages undergoes the next branching event after \(v\). The source reports that, under its exchangeability null hypothesis \(H_0\),
\[
\mathbb{P}_{H_0}(A_v)=\frac{m_v}{k_v}.
\]
This is a reported schematic probability relation, not a fully specified probability model. The counting convention for \(k_v\), the conditioning information, the meaning of exchangeability, and the treatment of simultaneous events, sampling, lineage termination, and absence of a subsequent branching event are unspecified. In particular, the source does not explain how the child lineages are identified among the lineages eligible for the next branching event. A precise interpretation requires these choices and a positive denominator. No binary-branching assumption is supplied.

The talk description says that these probabilities are used to construct a test statistic and test exchangeability, with illustrations using simulated data and SARS-CoV-2 sequence data. It supplies neither the statistic nor its null distribution, dependence assumptions, rejection rule, or a quantified conjecture. Thus no general implication, counterexample question, or unresolved theorem can be faithfully extracted from the heading alone.

A separate, truncated part of the record concerns estimating serial intervals using pathogen genomes and symptom-onset times while accounting for uncertainty about who infected whom. For an infecting individual \(A\) and an individual \(B\) infected by \(A\), with symptom-onset times \(s_A\) and \(s_B\), the described serial interval is
\[
I_{A,B}=s_B-s_A.
\]
The record supplies no sign restriction, statistical model, or complete mathematical target for this second topic; its final sentence is cut off.
\subsection{Short English statement}
The record describes C. Colijn's exchangeability test for timed phylogenies and reports the null probability \(m/k\) that a child of a branching node branches next. Essential model and testing details are missing. Although the supplied registry labels the entry open, its verification field says that no separate open mathematical problem was recovered. A second, truncated passage concerns estimating serial intervals from genomic and symptom-onset data.
\subsection{Sources}
\begin{itemize}
\item Local UnsolvedMath registry, permanent ID \(30005430\), source problem code \texttt{catalogue entry}, title \emph{Exchangeability Tests for Branching in Timed Phylogenies}. The supplied statement and original statement coincide; no clean statement is supplied.
\item \emph{Design and Analysis of Infectious Disease Studies} (2023), Oberwolfach Reports, DOI \texttt{10.4171/owr/2023/8}. The supplied citation states that the page range was not supplied by EMS. Original report: \url{https://doi.org/10.4171/owr/2023/8}. The first talk part, \emph{Does branching now imply branching next? Testing for exchangeability in timed phylogenies}, is attributed to C. Colijn as sole author. The second part, \emph{Estimating serial intervals with genomic data}, is attributed to a full authorship group whose names are absent from the supplied extract.
\item Supplied literature source: McNichol--Colijn, 2025 conference abstract, \url{https://dynamicsevolution.org/event/5/contributions/211/}.
\item Literature assessment dated exactly \texttt{2026-08-22}: the supplied curation labels the question open and describes the 2025 study as presenting a statistic and applications rather than a general affirmative or negative resolution. This is supplied dated curation, not an independently verified literature assessment. The background explicitly requests independent verification of time-sensitive status and cited claims.
\item The supplied statement-verification field separately identifies the question as a talk heading describing a developed test, rather than an unresolved mathematical problem, and records the statement status as \texttt{unrecoverable}. The original record is retained for provenance. No supplied citation establishes a precise standalone open conjecture.
\end{itemize}
\end{document}
